\documentclass[12pt,a4paper]{ctexbook}
\usepackage{geometry}
\geometry{margin=2.2cm}
\usepackage{setspace}
\setstretch{1.35}
\usepackage{hyperref}
\title{全宋词选·ci.song.6000}
\author{古典文献整理}
\date{}
\begin{document}
\maketitle
\tableofcontents
\newpage
\section{渔父词・渔父}
\textbf{周紫芝}\par\vspace{0.5em}
人间何物是穷通。\par
终向烟波作钓翁。\par
江不动，月横空。\par
漫郎船过小回中。\par

\section{花心动}
\textbf{李弥逊}\par\vspace{0.5em}
水馆风亭，晚香浓、一番荷芰经雨。\par
簟枕乍闲，襟裾初试，散尽满轩袢暑。\par
断云却送轻雷去，疏林外、玉钩微吐。\par
夜未阑，秋生败叶，暗催庭树。\par
天上佳期久阻。\par
星河畔，仙车缥缈云路。\par
旧恨未平，幽欢难驻，洒落半天风露。\par
绮罗人散金猊冷，醉魂到、华胥深处。\par
洞户悄，南楼画角自语。\par

\section{一寸金}
\textbf{李弥逊}\par\vspace{0.5em}
仙李盘根，自有云仍霭芳裔。\par
更留雨霜皮，临风玉树，紫髯丹颊，长生久视。\par
鹤帐琅书至。\par
长庚梦、当年暗记。\par
佳辰近，回首西风，渐喜秋英弄霜蕊。\par
暂卷双旌，鸣金吹竹，萱堂伴新戏。\par
对璧月流光，屏山供翠，碧云乍合，飞觞如缀。\par
早晚岩廊侍。\par
终不负、黄楼一醉。\par
丹青手、先与翻阶，万叶增春媚。\par

\section{蝶恋花}
\textbf{李弥逊}\par\vspace{0.5em}
百尺游丝当秀户。\par
不系春晖，只系闲愁住。\par
拾翠归来芳草路。\par
避人蝴蝶双飞去。\par
困脸羞眉无意绪。\par
陌上行人，记得清明否。\par
消息未来池阁暮。\par
濛濛一饷梨花雨。\par

\section{蝶恋花}
\textbf{李弥逊}\par\vspace{0.5em}
足力穷时山已晦。\par
却上轻舟，急棹穿沙背。\par
云影渐随风力退。\par
一川月白寒光碎。\par
唤客主人陶谢辈。\par
拂石移尊，不管游人醉。\par
罗绮丛中无此会。\par
只疑身在烟霞外。\par

\section{蝶恋花}
\textbf{李弥逊}\par\vspace{0.5em}
清晓天容争显晦。\par
溪上群山，戢戢分驼背。\par
谁似浮云知进退。\par
疏林嫩日黄金碎。\par
夜枕不眠憎鼠辈。\par
困眼贪晴，拚被风烟醉。\par
天意有情人不会。\par
分明置我风波外。\par

\section{蝶恋花}
\textbf{李弥逊}\par\vspace{0.5em}
百叠青山江一缕。\par
十里人家，路绕南台去。\par
榕叶满川飞白鹭。\par
疏帘半卷黄昏雨。\par
楼阁峥嵘天尺五。\par
荷芰风清，习习消袢暑。\par
老子人间无著处。\par
一尊来作横山主。\par

\section{蝶恋花}
\textbf{李弥逊}\par\vspace{0.5em}
小小芙蕖红半展。\par
占早争先，不奈腰肢软。\par
罗袜凌波娇欲颤。\par
向人如诉闺中怨。\par
把酒与君成眷恋。\par
约束新荷，四面低歌扇。\par
不放游人偷眼盼。\par
鸳鸯叶底潜窥见。\par

\section{虞美人}
\textbf{李弥逊}\par\vspace{0.5em}
上阳迟日千门锁。\par
花外流莺过。\par
一番春去又经秋。\par
惟有深宫明月、照人愁。\par
暗中白发随芳草。\par
却恨容颜好。\par
更无魂梦到昭阳。\par
肠断一双飞燕、在雕粱。\par

\section{虞美人}
\textbf{李弥逊}\par\vspace{0.5em}
海棠开后春谁主。\par
日日催花雨。\par
可怜新绿遍残枝。\par
不见香腮和粉、晕燕脂。\par
去年携手听金缕。\par
正是花飞处。\par
老来先自不禁愁。\par
这样愁来欺老、几时休。\par

\section{虞美人}
\textbf{李弥逊}\par\vspace{0.5em}
金泥捍拨春声碎。\par
恨入相思泪。\par
醉欺秋水绿云斜。\par
浑似梦中重到、阿环家。\par
主人著意留春住。\par
不醉无归去。\par
只愁银烛晓生寒。\par
明日落花飞絮、满长安。\par

\section{虞美人}
\textbf{李弥逊}\par\vspace{0.5em}
年年江上清秋节。\par
盏面分霜月。\par
不堪对月已伤离。\par
那更梅花开后、海棠时。\par
剑溪难驻仙游路。\par
直上云霄去。\par
藕花恰莫碍行舟。\par
要趁潮头八月、到扬州。\par

\section{虞美人}
\textbf{李弥逊}\par\vspace{0.5em}
方壶小有人谁到。\par
底事春知早。\par
使君和气粲如花。\par
更筑百花深处、驻春华。\par
云关远枕苕溪浅。\par
未放归怀展。\par
看残红紫绿阴斜。\par
鸾凤干霄却上、玉皇家。\par

\section{青玉案}
\textbf{李弥逊}\par\vspace{0.5em}
杨花尽做难拘管。\par
也解趁、飞红伴。\par
骢马无情人渐远。\par
沙平浅渡，雨湿孤村，何处长亭晚。\par
欲凭桃叶传春怨。\par
算不似、斜风倩双燕。\par
纵得书来春又换。\par
只将心事，分付眉尖，寂寞梨花院。\par

\section{菩萨蛮}
\textbf{李弥逊}\par\vspace{0.5em}
小山娇翠低歌扇。\par
雏莺学语春犹浅。\par
无力响方檀。\par
声随玉笋翻。\par
垂鬟云乍染。\par
媚靥香微点。\par
未解作轻颦。\par
凝情已动人。\par

\section{菩萨蛮}
\textbf{李弥逊}\par\vspace{0.5em}
凉飙轻散馀霞绮。\par
疏星冷浸明河水。\par
敧枕画檐风。\par
秋生草际蛩。\par
雁门离塞晚。\par
不道衡阳远。\par
归恨隔重山。\par
楼高莫凭栏。\par

\section{菩萨蛮}
\textbf{李弥逊}\par\vspace{0.5em}
风庭瑟瑟灯明灭。\par
碧梧枝上蝉声歇。\par
枕冷梦魂惊。\par
一阶寒水明。\par
鸟飞人未起。\par
月露清如洗。\par
无语听残更。\par
愁从两鬓生。\par

\section{菩萨蛮}
\textbf{李弥逊}\par\vspace{0.5em}
馀霞收尽寒烟绿。\par
江山一片团明玉。\par
敧枕画楼风。\par
愁生草际蛩。\par
金茎秋未老。\par
两鬓吴霜早。\par
忍负广寒期。\par
清尊对语谁。\par

\section{浣溪沙}
\textbf{李弥逊}\par\vspace{0.5em}
小小茅茨隐翠微。\par
桥平双手弄涟漪。\par
好风还动去年枝。\par
得雨疏梅肥欲展，人家次第有芳菲。\par
惜花恰莫探春迟。\par

\section{浣溪沙}
\textbf{李弥逊}\par\vspace{0.5em}
箫鼓哀吟乐楚臣。\par
牙樯锦缆簇江漘。\par
调高彩笔逞尖新。\par
海角逢时伤老大，莫辞卮酒话情亲。\par
与君同是异乡人。\par

\section{浣溪沙}
\textbf{李弥逊}\par\vspace{0.5em}
向日南枝不奈晴。\par
无风绛雪自飘零。\par
画楼更作断肠声。\par
小侧金荷迎落蕊，高烧银烛照残英。\par
生愁斜月酒初醒。\par

\section{临江仙}
\textbf{李弥逊}\par\vspace{0.5em}
枝上子规催去旆，柳条偏系离情。\par
片云留雨锁愁城。\par
不堪明月夜，寂寞照南荣。\par
莫作东山今日计，风雷已促鹏程。\par
功成来伴赤松行。\par
却寻鸿雁侣，尊酒会如星。\par

\section{临江仙}
\textbf{李弥逊}\par\vspace{0.5em}
燕去莺来昏又晓，劳生莫负心期。\par
菊花何必待开时。\par
十分浮玉蚁，一拍贯珠词。\par
少借笔端烟雨力，不须露染风披。\par
芳心微露定因谁。\par
风流今太傅，萧洒古东篱。\par

\section{临江仙}
\textbf{李弥逊}\par\vspace{0.5em}
多病渊明刚止酒，不禁秋蕊浮香。\par
饮船歌板已兼忘。\par
吴霜羞鬓改，无语对红妆。\par
小拈青枝撩鼻观，绝胜娇额涂黄。\par
独醒滋味怕新凉。\par
归来烛影乱，敧枕听更长。\par

\section{临江仙}
\textbf{李弥逊}\par\vspace{0.5em}
一片花飞春已减，那堪万点愁人。\par
可能春便负闲身。\par
细思愁不饮，却是自辜春。\par
且共一尊追落蕊，犹胜陌上成尘。\par
杯行到手莫辞频。\par
杏花须记取，曾与此翁邻。\par

\section{临江仙}
\textbf{李弥逊}\par\vspace{0.5em}
试问花枝馀几许，卷帘细雨随人。\par
风光犹恋苦吟身。\par
海棠浑怯冷，为我强留春。\par
细听惜花歌白雪，不知盏面生尘。\par
吹开吹谢漫惊频。\par
少陵真有味，爱酒觅南邻。\par

\section{醉花阴}
\textbf{李弥逊}\par\vspace{0.5em}
翠箔阴阴笼画阁。\par
昨夜东风恶。\par
香迳慢春泥，南陌东郊，惆怅妨行乐。\par
伤春比似年时觉。\par
潘鬓新来薄。\par
何处不禁愁，雨滴花腮，和泪胭脂落。\par

\section{醉花阴}
\textbf{李弥逊}\par\vspace{0.5em}
紫菊红萸开犯早。\par
独占秋光老。\par
酝造一般清，比著芝兰，犹自争多少。\par
霜刀翦叶呈纤巧。\par
手拈迎人笑。\par
云鬓一枝斜，小阁幽窗，是处都香了。\par

\section{清平乐}
\textbf{李弥逊}\par\vspace{0.5em}
烛花催晓。\par
醉玉颓春酒。\par
一骑东风消息到。\par
占得鳖头龙首。\par
长安去路骎骎。\par
明朝跃马芳阴。\par
应是花繁莺巧，东君著意琼林。\par

\section{清平乐}
\textbf{李弥逊}\par\vspace{0.5em}
一帘红雨。\par
飘荡谁家去。\par
门外垂杨千万缕。\par
不把东风留住。\par
旧巢燕子来迟。\par
故园绿暗残枝。\par
肠断画桥烟水，此情不许春知。\par

\section{浪淘沙}
\textbf{李弥逊}\par\vspace{0.5em}
把酒挽芳时。\par
醉袖淋漓。\par
多情楚客为秋悲。\par
未抵香飘红褪也，独绕空枝。\par
天女宝刀迟。\par
露染风披。\par
翠云叠叠拥铢衣。\par
知道筠溪春寂寞，来慰相思。\par

\section{浪淘沙}
\textbf{李弥逊}\par\vspace{0.5em}
乐事信难逢。\par
莫放匆匆。\par
飞红撩乱减春容。\par
临水不禁频送客，风袖龙锺。\par
小阁画堂东。\par
绮绣相重。\par
尊前谁唱夏云峰。\par
醒后欲寻溪上路，烟水无穷。\par

\section{谒金门}
\textbf{李弥逊}\par\vspace{0.5em}
春又老。\par
愁似落花难扫。\par
一醉一回才忘了。\par
醒来还满抱。\par
此恨欲凭谁道。\par
柳外数声啼鸟。\par
只恐春风吹不到。\par
断云连碧草。\par

\section{滴滴金}
\textbf{李弥逊}\par\vspace{0.5em}
广平未拚心如铁。\par
恨梅花、隔年别。\par
化工翦水斗春风，似南枝和月。\par
长鲸一饮宁论石。\par
想高歌、醉瑶席。\par
几时归去共尊罍，看寒花连陌。\par

\section{诉衷情}
\textbf{李弥逊}\par\vspace{0.5em}
小桃初破两三花。\par
深浅散馀霞。\par
东君也解人意，次第到山家。\par
临水岸，一枝斜。\par
照笼纱。\par
可怜何事，苦爱施朱，减尽容华。\par

\section{好事近}
\textbf{李弥逊}\par\vspace{0.5em}
春苑杂花芳，诗老剩夸梅格。\par
谁道武陵深处，便不如姑射。\par
莫分红浅与红深，点点是春色。\par
生怕一番风雨，半飘零江国。\par

\section{鹤冲天}
\textbf{李弥逊}\par\vspace{0.5em}
篘玉液，酿花光。\par
来趁北窗凉。\par
为君小摘蜀葵黄。\par
一似嗅枝香。\par
饮中仙，山中相。\par
也道十分宫样。\par
一般时候最宜尝。\par
竹院月侵床。\par

\section{天仙子}
\textbf{李弥逊}\par\vspace{0.5em}
飞盖追春春约伫。\par
繁杏枝头红未雨。\par
小楼翠幕不禁风，芳草路。\par
无尘处。\par
明月满庭人欲去。\par
一醉邻翁须记取。\par
见说新妆桃叶女。\par
明年却对此花时，留不住。\par
花前语。\par
总向似花人付与。\par

\section{清平乐}
\textbf{李弥逊}\par\vspace{0.5em}
断桥缺月。\par
点点枝头雪。\par
画角吹残声未歇。\par
早是一年春别。\par
寿阳弄粉成妆。\par
柔肠结结丁香。\par
可怕真梅轻妒，游蜂说与何妨。\par

\section{清平乐}
\textbf{李弥逊}\par\vspace{0.5em}
推愁何计。\par
车下忘乘坠。\par
日上南枝春有意。\par
已讶红酥如缀。\par
儿童缓整馀杯。\par
芒鞋午夜重来。\par
素面应憎月冷，真香不逐风回。\par

\section{清平乐}
\textbf{李弥逊}\par\vspace{0.5em}
长红小白。\par
何处寻梅格。\par
日日南山云满额。\par
暗老一分春色。\par
广平赋罢瑶香。\par
细□花蕊商量。\par
留得笔端雨露，后来收拾群芳。\par

\section{点绛唇}
\textbf{李弥逊}\par\vspace{0.5em}
翦翦疏花，托根宛在长松底。\par
蔓柯相倚。\par
便有凌霄志。\par
丹凤忽来，小队迎秋起。\par
留无计。\par
待公归侍。\par
重与分红翠。\par

\section{虞美人}
\textbf{李弥逊}\par\vspace{0.5em}
梨花院落溶溶雨。\par
弱柳低金缕。\par
画檐风露为谁明。\par
青翼来时试问、董双成。\par
去年春酒为眉寿。\par
花影浮金斗。\par
不须更觅老人星。\par
但愿一年一上、一千龄。\par

\section{醉花阴}
\textbf{李弥逊}\par\vspace{0.5em}
池面芙蕖红散绮。\par
鹊噪朱门喜。\par
环佩响天风，香霭杯盘，更约麻姑侍。\par
尘寰不隔蓬莱水。\par
束带岩廊戏。\par
瘦鹤与长松，且伴臞仙，久住人间世。\par

\section{醉花阴}
\textbf{李弥逊}\par\vspace{0.5em}
帘卷西风轻雨外。\par
揖数峰横翠。\par
楼上地行仙，压玉为醪，旋摘黄金蕊。\par
一觞一阕千秋岁。\par
不愿封侯贵。\par
长伴紫髯翁，踏月吹箫，笑咏云山里。\par

\section{感皇恩}
\textbf{李弥逊}\par\vspace{0.5em}
花院小回廊，庭萱成行。\par
水面红妆翠绡帐。\par
蓬莱云近，风露一番清旷。\par
星郎来碧落，长庚象。\par
素志未酬，丹心益壮。\par
且醉真珠小槽酿。\par
不须皓齿，拍手狂歌清唱。\par
一尊为寿庆，羲皇上。\par

\section{感皇恩}
\textbf{李弥逊}\par\vspace{0.5em}
密竹翦轻绡，华堂初建。\par
卷上虾须待开宴。\par
寿期春聚，芍药一番开遍。\par
砌成锦步帐，笼弦管。\par
绛节近颁，丹雏重见。\par
花里双双乍归燕。\par
重重乐事，凭仗东风拘管。\par
一时分付与，金荷劝。\par

\section{小重山}
\textbf{李弥逊}\par\vspace{0.5em}
鞭凤骖鸾自斗杓。\par
老君亲抱送，下层霄。\par
人间仙李占春饶。\par
千秋里，松月伴吹箫。\par
故国水云遥。\par
谪仙丹荔熟，剥红绡。\par
南山影转卧金蕉。\par
倾寒绿，眉寿比山高。\par

\section{小重山}
\textbf{李弥逊}\par\vspace{0.5em}
星斗心胸锦绣肠。\par
厌随尘土客，逐炎凉。\par
江山风月伴行藏。\par
无人识，高卧水云乡。\par
肘后有仙方。\par
假饶丹未就，寿须长。\par
儒冠多误莫思量。\par
十分酒，菡萏小池塘。\par

\section{花心动}
\textbf{李弥逊}\par\vspace{0.5em}
红日当楼，绣屏开、风裀舞花随步。\par
绛拂佩兰，香染妆梅，彷佛紫烟真侣。\par
雪消池馆年年会，玻璃泛、小槽新注。\par
弄箫语。\par
云璈未彻，暖回芳树。\par
瑶检曾攽寿缕。\par
好系日萦春，驭鸾深驻。\par
桂殿影寒，蓬山波阔，未似彩衣庭户。\par
坐看鹤鬓云来戏，重重拜、天阶雨露。\par
纵游处。\par
人间遍寻洞府。\par

\section{渔家傲}
\textbf{李弥逊}\par\vspace{0.5em}
海角秋高风力骤。\par
楼台四面山容瘦。\par
季子貂裘寒欲透。\par
悬弧昼。\par
高歌棣萼聊卮酒。\par
且共追欢宽白首。\par
清闲赢得身长久。\par
世上功名翻覆手。\par
为君寿。\par
腰间要看悬金斗。\par

\section{阮郎归}
\textbf{李弥逊}\par\vspace{0.5em}
黄花犹未拆霜枝。\par
今年秋较迟。\par
六么催泛玉东西。\par
登高庆诞时。\par
神仙事，古来稀。\par
且为千岁期。\par
戏莱堂上两庞眉。\par
何妨举案齐。\par

\section{醉落托}
\textbf{李弥逊}\par\vspace{0.5em}
霜林变绿。\par
画帘桂子排香粟。\par
一声檀板惊飞鹜。\par
弦管楼高，谁在阑干曲。\par
人生一笑难相属。\par
满堂何必堆金玉。\par
但求身健儿孙福。\par
鹤发年年，同泛清尊菊。\par

\section{柳梢青}
\textbf{李弥逊}\par\vspace{0.5em}
寿烟笼席。\par
采莲新按，舞腰无力。\par
占尽风光，人间天上，今夕何夕。\par
蓝袍换了莱衣，庆岁岁、君恩屡锡。\par
连夜欢声，满城佳气，和春留得。\par

\section{点绛唇}
\textbf{李弥逊}\par\vspace{0.5em}
花信争先，暗将春意传桃李。\par
寿卿同醉。\par
绿野连珠履。\par
麟阁丹青，眷注耆英裔。\par
眉间喜。\par
日边飞骑。\par
来促东山起。\par

\section{十样花}
\textbf{李弥逊}\par\vspace{0.5em}
陌上风光浓处。\par
第一寒梅先吐。\par
待得春来也，香销减，态凝伫。\par
百花休谩妒。\par
陌上风光浓处。\par
繁杏枝头春聚。\par
艳态最娇娆，堪比并，东邻女。\par
红梅何足数。\par
陌上风光浓处。\par
日暖山樱红露。\par
结子点朱唇，花谢后，君看取。\par
流莺偏嘱付。\par
陌上风光浓处。\par
忘却桃园归路。\par
洞口水流迟，香风动，红无数。\par
吹愁何处去。\par
陌上风光浓处。\par
最是海棠风措。\par
翠袖衬轻红，盈盈泪，怨春去。\par
黄昏微带雨。\par
陌上风光浓处。\par
自有花王为主。\par
富艳压群芳，蜂蝶戏，燕莺语。\par
东君都付与。\par
陌上风光浓处。\par
红药一番经雨。\par
把酒绕芳丛，花解语。\par
劝春住。\par
莫教容易去。\par

\section{菩萨蛮}
\textbf{李弥逊}\par\vspace{0.5em}
江城烽火连三月。\par
不堪对酒长亭别。\par
休作断肠声。\par
老来无泪倾。\par
风高帆影疾。\par
目送舟痕碧。\par
锦字几时来。\par
薰风无雁回。\par

\section{水调歌头}
\textbf{王以宁}\par\vspace{0.5em}
岁晚橘洲上，老叶舞愁红。\par
西山光翠，依旧影落酒杯中。\par
人在子亭高处，下望长沙城郭，猎猎酒帘风。\par
远水湛寒碧，独钓绿蓑翁。\par
怀往事，追昨梦，转头空。\par
孙郎前日，豪健颐指五都雄。\par
起拥奇才剑客，十万银戈赤帻，歌鼓壮军容。\par
何似裴相国，谈道老圭峰。\par

\section{水调歌头}
\textbf{王以宁}\par\vspace{0.5em}
大别我知友，突兀起西州。\par
十年重见，依旧秀色照清眸。\par
常记鲒碕狂客，邀我登楼雪霁，杖策拥羊裘。\par
山吐月千仞，残夜水明楼。\par
黄粱梦，未觉枕，几经秋。\par
与君邂逅，相逐飞步碧山头。\par
举酒一觞今古，叹息英雄骨冷，清泪不能收。\par
鹦鹉更谁赋，遗恨满芳洲。\par

\section{满庭芳}
\textbf{王以宁}\par\vspace{0.5em}
山耸方壶，潮通碧海，江东自昔名家。\par
玉真仙子，珰佩粲朝霞。\par
一种天香胜味，笑杨梅、不数枇杷。\par
难模写，牟尼妙质，光透紫丹砂。\par
咨嗟。\par
如此辈，不知何为，留滞天涯。\par
料甘心远引，无意纷华。\par
一任姚黄魏紫，供吟赏、银烛笼纱。\par
南游士，日餐千颗，不愿九霞车。\par

\section{满庭芳}
\textbf{王以宁}\par\vspace{0.5em}
千古南阳，刘郎乡国，依约楚俗秦风。\par
英姿豪气，耆旧笑谈中。\par
珰佩来从帝所，许洲花、潭菊从容。\par
霜秋晓，凉生日观，极目送飞鸿。\par
主公。\par
天下士，挥豪万字，一饮千锺。\par
醉高歌起舞，唤醒人龙。\par
我自人间漫浪，平生事、南北西东。\par
辞公去，寒眸激电，曾识小安丰。\par

\section{满庭芳}
\textbf{王以宁}\par\vspace{0.5em}
五十七年，侵寻老矣，小庵初筑林坰。\par
故人相过，喜雪舞祥霙。\par
遍野跳珠浅玉，纵儿童、收满金瓶。\par
明年待，洗光银海，袖手看升平。\par
先生。\par
齐物久，蚁丘罢战，蜗角休征。\par
趁尊前身健，有酒须倾。\par
随分村歌社舞，何须问、武宿文星。\par
忘怀矣，未能忘酒，相与醉忘形。\par

\section{满庭芳}
\textbf{王以宁}\par\vspace{0.5em}
千古黄州，雪堂奇胜，名与赤壁齐高。\par
竹楼千字，笔势压江涛。\par
笑问江头皓月，应曾照、今古英豪。\par
菖蒲酒，窊尊无恙，聊共访临皋。\par
陶陶。\par
谁晤对，粲花吐论，宫锦纫袍。\par
借银涛雪浪，一洗尘劳。\par
好在江山如画，人易老、双鬓难茠。\par
升平代，凭高望远，当赋反离骚。\par

\section{蓦山溪}
\textbf{王以宁}\par\vspace{0.5em}
平山堂上，侧盏歌南浦。\par
醉望五州山，渺千里、银涛东注。\par
钱郎英远，满腹贮精神，窥素壁，墨栖鸦，历历题诗处。\par
风裘雪帽，踏遍荆湘路。\par
回首古扬州，沁天外、残霞一缕。\par
德星光次，何日照长沙，渔夫曲，竹枝词，万古歌来暮。\par

\section{蓦山溪}
\textbf{王以宁}\par\vspace{0.5em}
雕弓绣帽。\par
戏马秦淮道。\par
风入马啼轻，曾踏遍、淮堤芳草。\par
飞英点点，春事已阑珊，风雨横，别离多，断送英雄老。\par
功名终在，休惜芳尊倒。\par
谈笑下燕云，看千里、风驱电扫。\par
男儿此事，莫待鬓丝棼，汉都护，万年觞，玉殿春风早。\par

\section{念奴娇}
\textbf{王以宁}\par\vspace{0.5em}
天工何意，碎琼珰玉佩，书空千尺。\par
箬笠蓑衫扁舟下，淮口烟林如织。\par
飞观嶙峋，子亭突兀，影浸澄淮碧。\par
纶巾鹤氅，是谁独笑携策。\par
遥想易水燕山，有人方醉赏，六花如席。\par
云重天低酣歌罢，胆壮乾坤犹窄。\par
射雉归来，铁鳞十万，踏碎千山白。\par
紫箫声断，唤回春满南陌。\par

\section{念奴娇}
\textbf{王以宁}\par\vspace{0.5em}
云收天碧。\par
渐风高露冷，群喧初寂。\par
远浦归舟荒渡口，搴篷横棹堤侧。\par
一带澄江，十分蟾影，千里寒光白。\par
襟怀迎爽，有人独步携策。\par
遥想帘卷琼楼，凭阑凝望，此意知何极。\par
声断琴箫双凤驾，凌厉飞仙同籍。\par
问汉乘槎，采珠临海，往事空踪迹。\par
莫辞吟赏，终宵清景难得。\par

\section{念奴娇}
\textbf{王以宁}\par\vspace{0.5em}
晚烟凝碧。\par
渐渔村山市，人归寂寂。\par
有客飞舟还顾访，应讶纶巾敧侧。\par
得意忘年，推诚投分，高论追元白。\par
英标逸气，笑予穷抱真策。\par
兴尽又复言归，秋风分袂，浩荡思无极。\par
咫尺昭山明翠壁，那知中隐咸籍。\par
说梦难听，闭门寻梦，肯念栖萍迹。\par
浪吟狂醉，几时还共重得。\par

\section{鹧鸪天}
\textbf{王以宁}\par\vspace{0.5em}
昔有书生荐寿杯。\par
清词妙绝贺方回。\par
黄花也欲为君寿，先向重阳六日开。\par
跻楚俗，上春台。\par
政成归去位三台。\par
明年寿酒君王劝，知有传宣敕使来。\par

\section{鹧鸪天}
\textbf{王以宁}\par\vspace{0.5em}
帝乙何年骑玉龙。\par
武夷仙伯笑相从。\par
长庚瑞应游仙梦，碧藕花开解愠风。\par
歌既醉，乐元丰。\par
明良相悦寿无穷。\par
野人更有深深祝，笑指三台十八公。\par

\section{鹧鸪天}
\textbf{王以宁}\par\vspace{0.5em}
桃李纷纷春事催。\par
桐花风定牡丹开。\par
天麟下作人间瑞，玉燕清宵入梦来。\par
红玉酎，紫霞杯。\par
五云深处望三台。\par
君家况有庭前柏，好个擎天八柱材。\par

\section{临江仙}
\textbf{王以宁}\par\vspace{0.5em}
眼看西园红与紫，年来几度芳菲。\par
吾生四十渐知非。\par
祗思青箬笠，江上雨霏霏。\par
饮酒但知寻夏季，不须远慕安期。\par
丹成仙去是何时。\par
龟藏何必学，春日正迟迟。\par

\section{临江仙}
\textbf{王以宁}\par\vspace{0.5em}
此理循环如引锯，春来百草菲菲。\par
道人一笑悟前非。\par
功名真长物，夕霭与朝霏。\par
梦褥清孙今绿隐，漫郎自许风期。\par
江楼景物得旬时。\par
平芜三百里，天阔夕阳迟。\par

\section{临江仙}
\textbf{王以宁}\par\vspace{0.5em}
闻道洛阳花正好，家家遮户春风。\par
道人饮处百壶空。\par
年年花下醉，看尽几番红。\par
此拐又从何处去，飘蓬一任西东。\par
语声虽异笑声同。\par
一轮清夜月，何处不相逢。\par

\section{浣溪沙}
\textbf{王以宁}\par\vspace{0.5em}
起看船头蜀锦张。\par
沙汀红叶舞斜阳。\par
杖拿惊起睡鸳鸯。\par
木落群山雕玉□，霜和冷月浸澄江。\par
疏篷今夜梦潇湘。\par

\section{西江月慢}
\textbf{吕渭老}\par\vspace{0.5em}
春风淡淡，清昼永、落英千尺。\par
桃杏散平郊，晴蜂来往，妙香飘掷。\par
傍画桥、煮酒青帘，绿杨风外，数声长笛。\par
记去年、紫陌朱门，花下旧相识。\par
向宝帕、裁书凭燕翼。\par
望翠阁、烟林似织。\par
闻道春衣犹未整，过禁烟寒食。\par
但记取、角枕情题，东窗休误，这些端的。\par
更莫待、青子绿阴春事寂。\par

\section{思佳客}
\textbf{吕渭老}\par\vspace{0.5em}
曾醉扬州十里楼。\par
竹西歌吹至今愁。\par
燕衔柳絮春心远，鱼入晴江水自流。\par
情渺渺，梦悠悠。\par
重寻罗带认银钩。\par
挂帆欲伴渔人去，只恐桃花误客舟。\par

\section{思佳客}
\textbf{吕渭老}\par\vspace{0.5em}
薄薄山云欲湿花。\par
双双燕子入帘斜。\par
西楼尚记垂垂雪，酌酒犹残片片霞。\par
人已远，发成华。\par
小楼疏竹宅谁家。\par
试凭去雁通消息，仙子当乘八月槎。\par

\section{夜游宫}
\textbf{吕渭老}\par\vspace{0.5em}
帘外繁霜未扫。\par
楼角动、玉绳横晓。\par
百和交焚瑞烟绕。\par
霁霞明，画屏深，天渺渺。\par
喜色连池沼。\par
荐眉寿、玉儿娇小。\par
早晚除书下天表。\par
日初长，莫等闲，孤一笑。\par

\section{浣溪沙}
\textbf{吕渭老}\par\vspace{0.5em}
风扫长林雪压枝。\par
纷纷冻鹊傍帘飞。\par
一尊聊作破寒威。\par
春意正愁梅漏泄，客情尤怕病禁持。\par
曲阑干外日初迟。\par

\end{document}